%% notes-tex/030-scaling-plots/sections/3-compute.tex

\section{Compute\secstatus{todo}}
\label{sec:compute}

The compute callback is what lets the grid say anything about $C$. It measures
FLOPs on five real training steps of every run, accumulates them per epoch with
the records seen and the GPU hours, and saves its state with the checkpoint, so
a chained point's totals continue. Figure~\ref{fig:compute} is what those keys
draw.

\tcfig{figures/030-scaling-plots-compute.pdf}{%
The compute family, sketched. \textbf{a}, C1: per-epoch validation loss against
records seen, one line per point colored by size, one panel per subset in the
real figure; the dot on each line is its selected epoch and the vertical line is
the matched budget $B$. \textbf{b}, C2: the same trajectories against cumulative
FLOPs; the dashed lower envelope is the compute-efficient frontier, the size
that reaches a given loss for the least compute. \textbf{c}, C3: the same against
GPU hours, the unit the queue charges; the horizontal offset of a size between b
and c is its utilization. \textbf{d}, C4: FLOPs per record and records per
second against size from the calibration jobs, with the spread over the five
measured steps as a whisker; a flat pace line across sizes means the loader, not
the GPU, sets the pace.
}{fig:compute}
\src{notes/assets/drawio/030-scaling-plots.gen.py}

\notesec{C1, learning curves against records seen}
The plot every other read is taken from. It shows where each point's selected
epoch sits, where $B$ cuts each curve, and which curves are still falling at
their last epoch. A full-pool curve that is still falling at ten epochs is the
picture of the unconverged headline point, and it is the plot to show beside any
exponent fit on the selected read.

\notesec{C2 and C3, the envelope}
Against FLOPs, the lower envelope of the trajectories across sizes is the
compute-efficient frontier of this architecture on this data: at each budget the
size whose curve is lowest. The grid has one budget line of an IsoFLOP sweep,
the fixed-epoch column, so the envelope here is read as the shape of the
frontier rather than its slope, which keeps C2 partial. Against GPU hours the
envelope is the same picture in the unit the queue charges, and because the FLOP
count excludes the loader while the wall clock includes it, the ratio of a size's
position in C2 to its position in C3 is the utilization that size achieved.

\notesec{C4, cost and pace against size}
From the two calibration jobs now, and from every grid run's first epoch after.
FLOPs per record should rise with $N$; whether records per second falls with it
is the question the calibration was queued to answer. A pace that follows the
FLOPs means the GPU binds and the small sizes could have trained longer in their
48 h; a pace that stays flat means the loader binds and no size finishes sooner
than 180x8. The whisker is the minimum to maximum over the five measured steps,
which on the per-entry path is the spread one batch's row count introduces.
